\documentclass[10pt,a4paper]{amsart}
\usepackage{amsmath,amssymb,amsthm,graphicx}
\usepackage{verbatim}
\newenvironment{reference}{\comment}{\endcomment}
%\newenvironment{reference}{}{}
\newenvironment{slogan}{\comment}{\endcomment}
\newenvironment{history}{\comment}{\endcomment}

% For commutative diagrams we use Xy-pic
\usepackage[all]{xy}

% We use 2cell for 2-commutative diagrams.
\xyoption{2cell}
% Theorem environments.
%
\theoremstyle{plain}
\newtheorem{theorem}[subsection]{Theorem}
\newtheorem{proposition}[subsection]{Proposition}
\newtheorem{lemma}[subsection]{Lemma}

\theoremstyle{definition}
\newtheorem{definition}[subsection]{Definition}
\newtheorem{example}[subsection]{Example}
\newtheorem{exercise}[subsection]{Exercise}
\newtheorem{situation}[subsection]{Situation}

\theoremstyle{remark}
\newtheorem{remark}[subsection]{Remark}
\newtheorem{remarks}[subsection]{Remarks}

\numberwithin{equation}{subsection}

% Macros
%
\def\lim{\mathop{\mathrm{lim}}\nolimits}
\def\colim{\mathop{\mathrm{colim}}\nolimits}
\def\Spec{\mathop{\mathrm{Spec}}}
\def\Hom{\mathop{\mathrm{Hom}}\nolimits}
\def\Ext{\mathop{\mathrm{Ext}}\nolimits}
\def\SheafHom{\mathop{\mathcal{H}\!\mathit{om}}\nolimits}
\def\SheafExt{\mathop{\mathcal{E}\!\mathit{xt}}\nolimits}
\def\Sch{\mathit{Sch}}
\def\Mor{\mathop{\mathrm{Mor}}\nolimits}
\def\Ob{\mathop{\mathrm{Ob}}\nolimits}
\def\Sh{\mathop{\mathit{Sh}}\nolimits}
\def\NL{\mathop{N\!L}\nolimits}
\def\CH{\mathop{\mathrm{CH}}\nolimits}
\def\proetale{{pro\text{-}\acute{e}tale}}
\def\etale{{\acute{e}tale}}
\def\QCoh{\mathit{QCoh}}
\def\Ker{\mathop{\mathrm{Ker}}}
\def\Im{\mathop{\mathrm{Im}}}
\def\Coker{\mathop{\mathrm{Coker}}}
\def\Coim{\mathop{\mathrm{Coim}}}

% Boxtimes
%
\DeclareMathSymbol{\boxtimes}{\mathbin}{AMSa}{"02}

%
% Macros for moduli stacks/spaces
%
\def\QCohstack{\mathcal{QC}\!\mathit{oh}}
\def\Cohstack{\mathcal{C}\!\mathit{oh}}
\def\Spacesstack{\mathcal{S}\!\mathit{paces}}
\def\Quotfunctor{\mathrm{Quot}}
\def\Hilbfunctor{\mathrm{Hilb}}
\def\Curvesstack{\mathcal{C}\!\mathit{urves}}
\def\Polarizedstack{\mathcal{P}\!\mathit{olarized}}
\def\Complexesstack{\mathcal{C}\!\mathit{omplexes}}
% \Pic is the operator that assigns to X its picard group, usage \Pic(X)
% \Picardstack_{X/B} denotes the Picard stack of X over B
% \Picardfunctor_{X/B} denotes the Picard functor of X over B
\def\Pic{\mathop{\mathrm{Pic}}\nolimits}
\def\Picardstack{\mathcal{P}\!\mathit{ic}}
\def\Picardfunctor{\mathrm{Pic}}
\def\Deformationcategory{\mathcal{D}\!\mathit{ef}}
\usepackage{cleveref}
\crefname{originallemma}{original lemma}{original lemmas}
\Crefname{originallemma}{Original lemma}{Original lemmas}
\makeatletter
\newlabel{algebra-lemma-completion-flat}{{00MB}{source}{}{sourcepointer-1}{}}
\newlabel{algebra-lemma-completion-flat@cref}{{[originallemma][0][]00MB}{[1][0][]source}{}{}{}}
\makeatother
\crefname{originallemma}{original lemma}{original lemmas}
\Crefname{originallemma}{Original lemma}{Original lemmas}
\makeatletter
\newlabel{formal-spaces-lemma-completed-tensor-product}{{0GB4}{source}{}{sourcepointer-2}{}}
\newlabel{formal-spaces-lemma-completed-tensor-product@cref}{{[originallemma][0][]0GB4}{[1][0][]source}{}{}{}}
\makeatother
\crefname{originalremark}{original remark}{original remarks}
\Crefname{originalremark}{Original remark}{Original remarks}
\makeatletter
\newlabel{formal-spaces-remark-base-change-variant-Noetherian}{{0GBF}{source}{}{sourcepointer-3}{}}
\newlabel{formal-spaces-remark-base-change-variant-Noetherian@cref}{{[originalremark][0][]0GBF}{[1][0][]source}{}{}{}}
\makeatother
\crefname{originalremark}{original remark}{original remarks}
\Crefname{originalremark}{Original remark}{Original remarks}
\makeatletter
\newlabel{formal-spaces-remark-base-change-variant-variant-Noetherian}{{0GBG}{source}{}{sourcepointer-4}{}}
\newlabel{formal-spaces-remark-base-change-variant-variant-Noetherian@cref}{{[originalremark][0][]0GBG}{[1][0][]source}{}{}{}}
\makeatother
\crefname{originalreference}{original reference}{original references}
\Crefname{originalreference}{Original reference}{Original references}
\makeatletter
\newlabel{lemma-Noetherian-finite-type-red}{{0GC6}{source}{}{sourcepointer-5}{}}
\newlabel{lemma-Noetherian-finite-type-red@cref}{{[originalreference][0][]0GC6}{[1][0][]source}{}{}{}}
\makeatother
\crefname{originalreference}{original reference}{original references}
\Crefname{originalreference}{Original reference}{Original references}
\makeatletter
\newlabel{lemma-finite-type}{{04U5}{source}{}{sourcepointer-6}{}}
\newlabel{lemma-finite-type@cref}{{[originalreference][0][]04U5}{[1][0][]source}{}{}{}}
\makeatother
\crefname{originalreference}{original reference}{original references}
\Crefname{originalreference}{Original reference}{Original references}
\makeatletter
\newlabel{lemma-finite-type-finite-type-red}{{0GC5}{source}{}{sourcepointer-7}{}}
\newlabel{lemma-finite-type-finite-type-red@cref}{{[originalreference][0][]0GC5}{[1][0][]source}{}{}{}}
\makeatother
\crefname{originalreference}{original reference}{original references}
\Crefname{originalreference}{Original reference}{Original references}
\makeatletter
\newlabel{lemma-finite-type-red}{{0GC2}{source}{}{sourcepointer-8}{}}
\newlabel{lemma-finite-type-red@cref}{{[originalreference][0][]0GC2}{[1][0][]source}{}{}{}}
\makeatother
\crefname{originalreference}{original reference}{original references}
\Crefname{originalreference}{Original reference}{Original references}
\makeatletter
\newlabel{lemma-flat-axioms}{{0GC9}{source}{}{sourcepointer-9}{}}
\newlabel{lemma-flat-axioms@cref}{{[originalreference][0][]0GC9}{[1][0][]source}{}{}{}}
\makeatother
\crefname{originallemma}{original lemma}{original lemmas}
\Crefname{originallemma}{Original lemma}{Original lemmas}
\makeatletter
\newlabel{more-algebra-lemma-limit-flat}{{0912}{source}{}{sourcepointer-10}{}}
\newlabel{more-algebra-lemma-limit-flat@cref}{{[originallemma][0][]0912}{[1][0][]source}{}{}{}}
\makeatother
\title{The flatness property P is stable under base change by the finite type properties Q and P+R}
\author{The Stacks Project Authors}
\date{stacks-project:restricted.tex:local-snapshot-20261012; GNU FDL 1.2 or later}
\begin{document}
\maketitle
\noindent\textit{Editorial note.} All delivered mathematical text is the author's own, copied from the pinned Stacks Project chapter file restricted.tex; no statement or proof step was rewritten, reordered or removed.\par\medskip
\noindent\textit{Editorial note.} The counted proof is the author's complete argument: the implication R implies Q is used as a black-box in-source citation (lemma-finite-type-finite-type-red); this card does not depend on that lemma's own proof\par\medskip
\noindent\textit{Editorial note.} The label lemma-base-change-flat-continuous names the counted statement, and the proof environment immediately following it in the source is delivered verbatim as its proof.\par\medskip
\noindent\textit{Editorial note.} The counted statement is carried by the author's own lemma environment in the chapter source.\par\medskip
\noindent\textit{Allowed standard prerequisites (editorial).}\begin{itemize}
\item Lemma lemma-flat-axioms (in-source definition carrier): defines property P (flatness of the relevant localizations) on WAdm\textasciicircum{}Noeth
\item Lemma lemma-finite-type-red and lemma-finite-type (in-source definition carriers): define properties Q and R
\item Formal Spaces, Lemma formal-spaces-lemma-completed-tensor-product: Noetherianity of the completed tensor under Q
\item Lemma lemma-Noetherian-finite-type-red (in-source): Q implies topologically finite type
\item Algebra, Lemma algebra-lemma-completion-flat: the completion of a finite type map is flat
\item standard completed-tensor quotient identity: A (tensor)\_B C/J\textasciicircum{}n = (A (completed tensor)\_B C)/J\textasciicircum{}n(...)
\item a limit lemma: a J-adically complete module whose quotients M/J\textasciicircum{}n M are flat over C/J\textasciicircum{}n is flat over C
\end{itemize}\medskip
\ifdefined\hypertarget\hypertarget{sourcepointer-1}{}\fi
\noindent\textit{Source reference.} Original lemma 00MB: sources/stacks-chow/tags.txt, source lines 817--817. Prerequisite pointer invoked at 2992; it names a standard fact of a sibling chapter that the counted argument uses, and every step of the counted proof built on it is retained verbatim in the delivered body.\par\medskip
\ifdefined\hypertarget\hypertarget{sourcepointer-2}{}\fi
\noindent\textit{Source reference.} Original lemma 0GB4: sources/stacks-chow/tags.txt, source lines 20359--20359. Prerequisite pointer invoked at 2974; it names a standard fact of a sibling chapter that the counted argument uses, and every step of the counted proof built on it is retained verbatim in the delivered body.\par\medskip
\ifdefined\hypertarget\hypertarget{sourcepointer-3}{}\fi
\noindent\textit{Source reference.} Original remark 0GBF: sources/stacks-chow/tags.txt, source lines 20370--20370. Prerequisite pointer invoked at 2964; it names a standard fact of a sibling chapter that the counted argument uses, and every step of the counted proof built on it is retained verbatim in the delivered body.\par\medskip
\ifdefined\hypertarget\hypertarget{sourcepointer-4}{}\fi
\noindent\textit{Source reference.} Original remark 0GBG: sources/stacks-chow/tags.txt, source lines 20371--20371. Prerequisite pointer invoked at 2961; it names a standard fact of a sibling chapter that the counted argument uses, and every step of the counted proof built on it is retained verbatim in the delivered body.\par\medskip
\ifdefined\hypertarget\hypertarget{sourcepointer-5}{}\fi
\noindent\textit{Source reference.} Original reference 0GC6: sources/stacks-restricted/pinned-original-source.tex, source lines 2749--2749. Prerequisite pointer invoked at 2982; it names a standard fact of the same chapter that the counted argument uses, and every step of the counted proof built on it is retained verbatim in the delivered body.\par\medskip
\ifdefined\hypertarget\hypertarget{sourcepointer-6}{}\fi
\noindent\textit{Source reference.} Original reference 04U5: sources/stacks-restricted/pinned-original-source.tex, source lines 2439--2439. Prerequisite pointer invoked at 2956; it names a standard fact of the same chapter that the counted argument uses, and every step of the counted proof built on it is retained verbatim in the delivered body.\par\medskip
\ifdefined\hypertarget\hypertarget{sourcepointer-7}{}\fi
\noindent\textit{Source reference.} Original reference 0GC5: sources/stacks-restricted/pinned-original-source.tex, source lines 2738--2738. Prerequisite pointer invoked at 2975; it names a standard fact of the same chapter that the counted argument uses, and every step of the counted proof built on it is retained verbatim in the delivered body.\par\medskip
\ifdefined\hypertarget\hypertarget{sourcepointer-8}{}\fi
\noindent\textit{Source reference.} Original reference 0GC2: sources/stacks-restricted/pinned-original-source.tex, source lines 2631--2631. Prerequisite pointer invoked at 2954; it names a standard fact of the same chapter that the counted argument uses, and every step of the counted proof built on it is retained verbatim in the delivered body.\par\medskip
\ifdefined\hypertarget\hypertarget{sourcepointer-9}{}\fi
\noindent\textit{Source reference.} Original reference 0GC9: sources/stacks-restricted/pinned-original-source.tex, source lines 2890--2890. Prerequisite pointer invoked at 2953; it names a standard fact of the same chapter that the counted argument uses, and every step of the counted proof built on it is retained verbatim in the delivered body.\par\medskip
\ifdefined\hypertarget\hypertarget{sourcepointer-10}{}\fi
\noindent\textit{Source reference.} Original lemma 0912: sources/stacks-chow/tags.txt, source lines 11430--11430. Prerequisite pointer invoked at 2999; it names a standard fact of a sibling chapter that the counted argument uses, and every step of the counted proof built on it is retained verbatim in the delivered body.\par\medskip
\begin{lemma}
\label{lemma-base-change-flat-continuous}
Denote $P$ the property of arrows of $\textit{WAdm}^{Noeth}$
defined in Lemma \ref{lemma-flat-axioms}.
Denote $Q$ the property defined in Lemma \ref{lemma-finite-type-red}
viewed as a property of arrows of $\textit{WAdm}^{Noeth}$.
Denote $R$ the property defined in Lemma \ref{lemma-finite-type}
viewed as a property of arrows of $\textit{WAdm}^{Noeth}$. Then
\begin{enumerate}
\item $P$ is stable under base change by $Q$
(Formal Spaces, Remark
\ref{formal-spaces-remark-base-change-variant-variant-Noetherian}), and
\item $P + R$ is stable under base change
(Formal Spaces, Remark
\ref{formal-spaces-remark-base-change-variant-Noetherian}).
\end{enumerate}
\end{lemma}

\begin{proof}
The statement makes sense as each of the properties $P$, $Q$, and $R$
is a local property of morphisms of $\textit{WAdm}^{Noeth}$.
Let $\varphi : B \to A$ and $\psi : B \to C$ be morphisms of
$\textit{WAdm}^{Noeth}$. If either $Q(\varphi)$ or $Q(\psi)$
then we see that $A \widehat{\otimes}_B C$ is Noetherian by
Formal Spaces, Lemma \ref{formal-spaces-lemma-completed-tensor-product}.
Since $R$ implies $Q$ (Lemma \ref{lemma-finite-type-finite-type-red}),
we find that this holds in both cases (1) and (2).
This is the first thing we have to check. It remains to show
that $C \to A \widehat{\otimes}_B C$ is flat.

\medskip\noindent
Proof of (1). Fix ideals of definition $I \subset A$ and $J \subset B$.
By Lemma \ref{lemma-Noetherian-finite-type-red} the ring map
$B \to C$ is topologically of finite type. Hence
$B \to C/J^n$ is of finite type for all $n \geq 1$.
Hence $A \otimes_B C/J^n$ is Noetherian as a ring
(because it is of finite type over $A$ and $A$ is Noetherian).
Thus the $I$-adic completion $A \widehat{\otimes}_B C/J^n$
of $A \otimes_B C/J^n$ is flat over $C/J^n$ because
$C/J^n \to A \otimes_B C/J^n$ is flat as a base change of
$B \to A$ and because
$A \otimes_B C/J^n \to A \widehat{\otimes}_B C/J^n$ is flat by
Algebra, Lemma \ref{algebra-lemma-completion-flat}
Observe that $A \widehat{\otimes}_B C/J^n =
(A \widehat{\otimes}_B C)/J^n(A \widehat{\otimes}_B C)$; details omitted.
We conclude that $M = A \widehat{\otimes}_B C$ is a $C$-module
which is complete with respect to the $J$-adic topology
such that $M/J^nM$ is flat over $C/J^n$ for all $n \geq 1$.
This implies that $M$ is flat over $C$ by
More on Algebra, Lemma \ref{more-algebra-lemma-limit-flat}.

\medskip\noindent
Proof of (2). In this case $B \to A$ is adic and hence we have just
$A \widehat{\otimes}_B C = \lim A \otimes_B C/J^n$.
The rings $A \otimes_B C/J^n$ are Noetherian by an application of 
Formal Spaces, Lemma \ref{formal-spaces-lemma-completed-tensor-product}
with $C$ replaced by $C/J^n$.
We conclude in the same manner as before.
\end{proof}

\end{document}
